% ============================================================
%  Příprava na maturitu – Mechatronika
%  Kompiluj: pdflatex maturita_mechatronika.tex
%  Kódování: UTF-8
% ============================================================

\documentclass[a4paper,10pt]{article}

% --- Balíčky ---
\usepackage[T1]{fontenc}
\usepackage[utf8]{inputenc}
\usepackage[czech]{babel}
\usepackage[left=2.4cm, right=1.8cm, top=1.8cm, bottom=1.8cm]{geometry}
\usepackage{lmodern}
\usepackage{microtype}
\usepackage{parskip}
\usepackage{titlesec}
\usepackage{enumitem}
\usepackage{xcolor}
\usepackage{hyperref}
\usepackage{tabularx}
\usepackage{booktabs}
\usepackage{array}
\usepackage{tikz}
\usepackage{eso-pic}
\usepackage{mdframed}
\usepackage{amsmath}
\usepackage{amssymb}
\usepackage{pgfplots}
\usepackage{karnaugh-map}
\usepackage{circuitikz}
\pgfplotsset{compat=1.18}
\usetikzlibrary{arrows.meta, positioning, shapes.geometric, calc}

% --- Barvy ---
\definecolor{accent}{HTML}{03254E}
\definecolor{stripeblue}{HTML}{568EA3}
\definecolor{medgray}{HTML}{888888}
\definecolor{lightblue}{HTML}{EAF2F8}
\definecolor{lightyellow}{HTML}{FDFBE4}
\definecolor{lightgreen}{HTML}{EAF7EA}

% --- Modrý pruh vlevo ---
\AddToShipoutPictureBG{%
  \begin{tikzpicture}[remember picture, overlay]
    \fill[stripeblue]
      (current page.north west)
      rectangle
      ([xshift=10mm] current page.south west);
    \fill[accent!60]
      ([xshift=10mm] current page.north west)
      rectangle
      ([xshift=12mm] current page.south west);
  \end{tikzpicture}%
}

\hypersetup{colorlinks=true, urlcolor=accent, linkcolor=accent}

\titleformat{\section}
  {\large\bfseries\color{accent}}
  {\thesection.}
  {0.5em}
  {}
  [\color{accent}\titlerule]

\titlespacing{\section}{0pt}{16pt}{6pt}

\newmdenv[
  backgroundcolor=lightblue,
  linecolor=stripeblue,
  linewidth=1pt,
  roundcorner=4pt,
  innerleftmargin=8pt,
  innerrightmargin=8pt,
  innertopmargin=6pt,
  innerbottommargin=6pt,
  skipabove=4pt,
  skipbelow=4pt
]{pojmybox}

\newmdenv[
  backgroundcolor=lightyellow,
  linecolor=accent!40,
  linewidth=1pt,
  roundcorner=4pt,
  innerleftmargin=8pt,
  innerrightmargin=8pt,
  innertopmargin=6pt,
  innerbottommargin=6pt,
  skipabove=4pt,
  skipbelow=8pt
]{vysvetlenibox}

\newmdenv[
  backgroundcolor=lightgreen,
  linecolor=accent!30,
  linewidth=1pt,
  roundcorner=4pt,
  innerleftmargin=8pt,
  innerrightmargin=8pt,
  innertopmargin=6pt,
  innerbottommargin=6pt,
  skipabove=4pt,
  skipbelow=8pt
]{vzorcebox}

\pagestyle{plain}

% ============================================================
\begin{document}

% --- Záhlaví ---
\begin{minipage}[t]{0.75\textwidth}
    {\Huge\bfseries\color{accent} Příprava na maturitu}\\[4pt]
    {\large\color{stripeblue} Mechatronika}\\[2pt]
    \textcolor{medgray}{\small Suchánek Lukáš \quad SPŠ Prosek \quad 2026}
\end{minipage}
\hfill
\begin{minipage}[t]{0.22\textwidth}
    \raggedleft
    \begin{tikzpicture}
        \draw[color=stripeblue, rounded corners=4pt, line width=1pt]
            (0,0) rectangle (3,1.2)
            node[midway, align=center, color=accent, font=\small\bfseries]
            {Otázek: 22};
    \end{tikzpicture}
\end{minipage}

\vspace{6pt}
\noindent\color{accent}\rule{\linewidth}{1.5pt}
\vspace{2pt}

% ============================================================
\section{Logické řízení -- programovací jazyky, logické funkce}

\begin{pojmybox}
\textbf{Klíčové pojmy:}
Booleova algebra, Karnaughova mapa, logické funkce (AND, OR, NOT, NAND, NOR, XOR),
pravdivostní tabulka, funkčně úplný systém
\end{pojmybox}

\begin{vysvetlenibox}
	\textbf{Úvod do logiky}
	\begin{itemize}[noitemsep]
		\item z řeckého slova logos $\Rightarrow$ důvod
		\item Věda zkoumající způsob vyvozování závěrů a způsob předpokladu
		\item Myšlenková cesta, která vede k daným závěrům
	\end{itemize}
	

	\textbf{Výrok}
	\begin{itemize}[noitemsep]
		\item Tvrzení u kterého má smysl ptát se, zda je \textbf{pravdivé} či \textbf{nepravdivé}
		\item Jednoduchý výrok $\rightarrow$ "Včera tady pršelo"
		\item Složený výrok $\rightarrow$ "Včera tady pršelo a současně se jmenuji Pepa"
	\end{itemize}
	
\textbf{Logické spojky}

\begin{center}
	\resizebox{\textwidth}{!}{
		\begin{tabular}{|l|l|p{4cm}|p{5.5cm}|}
			\hline
			\textbf{Operace} & \textbf{Zápis} & \textbf{Význam} & \textbf{Příklad} \\
			\hline
			Negace & $\neg A$, NOT A & není pravda A & $\neg(\text{Prší}) = \text{Neprší}$ \\
			\hline
			Konjunkce & $A \land B$, $A \cdot B$, A AND B & A a současně B & Nemám vlasy a současně se jmenuji Pepa \\
			\hline
			Disjunkce & $A \lor B$, A + B, A OR B & A nebo B & Vlak přijel na 1. nebo 2. kolej \\
			\hline
			Implikace & $A \Rightarrow B$ & jestliže A, pak B & Jestliže pršelo, je mokro \\
			\hline
			Ekvivalence & $A \Leftrightarrow B$ & A tehdy a jen tehdy B & Číslo je sudé tehdy a jen tehdy, pokud je dělitelné 2 \\
			\hline
		\end{tabular}
	}
\end{center}

\textbf{Důležité pojmy}
\begin{itemize}[noitemsep]
	\item funkce, definiční obor a obor hodnot
	\item logická proměnná
	\item hradlo
	\item logický obvod
\end{itemize}
	
\textbf{Booleova algebra}\\
Logické řízení pracuje s binárními signály: \textbf{0} (nepravda, vypnuto)
a \textbf{1} (pravda, zapnuto). Základem je Booleova algebra -- algebraická
struktura s operacemi součin (AND), součet (OR) a negace (NOT).

\textbf{Fyzikální reprezentace logických úrovní}
\begin{itemize}[noitemsep]
    \item \textbf{TTL logika} (Transistor-Transistor Logic):
          $0 = 0\text{--}0{,}8\,\text{V}$, $1 = 2\text{--}5\,\text{V}$
    \item \textbf{CMOS logika} (Complementary MOS):
          $0 = 0\text{--}1{,}5\,\text{V}$, $1 = 3{,}5\text{--}5\,\text{V}$
    \item \textbf{PLC vstupy} (podle normy IEC 61131-2):
          $0 = 0\text{--}5\,\text{V}$, $1 = 15\text{--}30\,\text{V}$
\end{itemize}

\textbf{Základní logické funkce a jejich pravdivostní tabulky}

\begin{center}
\begin{tabular}{cc|c|c|c|c|c}
    \toprule
    $A$ & $B$ & AND & OR & NAND & NOR & XOR \\
    \midrule
    0 & 0 & 0 & 0 & 1 & 1 & 0 \\
    0 & 1 & 0 & 1 & 1 & 0 & 1 \\
    1 & 0 & 0 & 1 & 1 & 0 & 1 \\
    1 & 1 & 1 & 1 & 0 & 0 & 0 \\
    \bottomrule
\end{tabular}
\end{center}

\textbf{Vysvětlení logických funkcí:}
\begin{itemize}[noitemsep, leftmargin=1.4em, topsep=4pt]
    \item \textbf{AND} (konjunkce, $\cdot$) -- výstup je 1, pouze když \textit{oba} vstupy jsou 1
    \item \textbf{OR} (disjunkce, $+$) -- výstup je 1, když \textit{alespoň jeden} vstup je 1
    \item \textbf{NOT} (negace, $\overline{A}$) -- převrátí hodnotu (0$\to$1, 1$\to$0)
    \item \textbf{NAND} (NOT-AND) -- negovaný AND; univerzální hradlo
    \item \textbf{NOR} (NOT-OR) -- negovaný OR; také univerzální
    \item \textbf{XOR} (exklusivní OR) -- výstup je 1, když se vstupy \textit{liší}
\end{itemize}

\textbf{Základní zákony Booleovy algebry}

\begin{vzorcebox}
\begin{align*}
A \cdot 0 &= 0 & A + 0 &= A & \text{(nulový prvek)} \\
A \cdot 1 &= A & A + 1 &= 1 & \text{(jednotkový prvek)} \\
A \cdot A &= A & A + A &= A & \text{(idempotence)} \\
A \cdot \overline{A} &= 0 & A + \overline{A} &= 1 & \text{(komplementarita)} \\
\overline{\overline{A}} &= A & & & \text{(dvojí negace)} \\
\overline{A \cdot B} &= \overline{A} + \overline{B} &
\overline{A + B} &= \overline{A} \cdot \overline{B} & \text{(De Morganovy zákony)}
\end{align*}
\end{vzorcebox}

\textbf{Funkčně úplné systémy}\\
NAND nebo NOR samy o sobě dokáží realizovat libovolnou logickou funkci.
To znamená, že kdybychom měli k dispozici pouze hradla NAND (nebo pouze NOR),
dokážeme sestavit jakýkoliv logický obvod.

\begin{center}
\begin{tikzpicture}[scale=0.8,
  gate/.style={rectangle, draw=accent, fill=lightblue, rounded corners=2pt,
               minimum width=1.2cm, minimum height=0.5cm, font=\small}]
  % NAND jako NOT
  \node[gate] (n1) at (0,0) {NAND};
  \draw[>-] (-0.8,-0.1) -- (n1.west);
  \draw[>-] (-0.8,0.1) -- (n1.west);
  \draw[>-] (n1.east) -- (0.8,0);
  \node[below, font=\scriptsize] at (0,-0.4) {NOT z NAND};

  % NAND jako AND
  \node[gate] (n2) at (2.5,0) {NAND};
  \node[gate] (n3) at (4.5,0) {NAND};
  \draw[>-] (1.7,-0.1) -- (n2.west);
  \draw[>-] (1.7,0.1) -- (n2.west);
  \draw[>-] (n2.east) -- (n3.west);
  \draw[>-] (n3.east) -- (n3.west);
  \draw[>-] (n3.east) -- (5.5,0);
  \node[below, font=\scriptsize] at (3.5,-0.4) {AND z NAND};
\end{tikzpicture}
\end{center}

\textbf{Programovací jazyky PLC dle normy IEC 61131-3}\\
Norma IEC 61131-3 definuje 5 programovacích jazyků pro PLC:

\begin{center}
\begin{tabularx}{0.95\linewidth}{@{} l l X @{}}
    \toprule
    \textbf{Zkratka} & \textbf{Název} & \textbf{Popis a použití} \\
    \midrule
    LD  & Ladder Diagram & Grafický, podobá se elektrickému schématu;
                           nejrozšířenější v průmyslu \\
    FBD & Function Block Diagram & Grafický, propojení funkčních bloků;
                                   vhodný pro regulaci \\
    ST  & Structured Text & Textový, podobný Pascalu;
                            vhodný pro složité algoritmy \\
    IL  & Instruction List & Nízkoúrovňový assembler PLC;
                             dnes již málo používaný \\
    SFC & Sequential Function Chart & Šablona stavového automatu;
                                       sekvenční procesy \\
    \bottomrule
\end{tabularx}
\end{center}

\textbf{Příklad kódu v jednotlivých jazycích:}

\textit{LD (Ladder Diagram):}
\begin{verbatim}
   I0.0      I0.1      Q0.0
|----| |------|/|------( )---|
|  start     stop     motor  |
\end{verbatim}

\textit{ST (Structured Text):}
\begin{verbatim}
IF start AND NOT stop THEN
    motor := TRUE;
ELSE
    motor := FALSE;
END_IF;
\end{verbatim}

\textit{SFC (Sequential Function Chart):}
\begin{verbatim}
[INIT] --(start)--> [RUN] --(stop)--> [STOP]
   ^                                   |
   +-----------------------------------+
\end{verbatim}

\textbf{Proč více jazyků?}\\
Různí výrobci PLC (Siemens, Allen-Bradley, Schneider) mají vlastní tradice.
Norma IEC 61131-3 umožňuje přenositelnost programů mezi platformami a usnadňuje
komunikaci mezi dodavateli.
\end{vysvetlenibox}

\newpage
\begin{vysvetlenibox}
	\textbf{Příklad}\\
	Jsou tři větráky a, b, c a alarm x, pokud dva a více větráků nefungují, spustí se alarm
	
	\textbf{Pravdivostní tabulka}
	\[
	\begin{array}{|c|c|c||c|}
		\hline
		A & B & C & X \\
		\hline
		0 & 0 & 0 & 1 \\
		0 & 0 & 1 & 1 \\
		0 & 1 & 0 & 1 \\
		0 & 1 & 1 & 0 \\
		\hline
		1 & 0 & 0 & 1 \\
		1 & 0 & 1 & 0 \\
		1 & 1 & 0 & 0 \\
		1 & 1 & 1 & 0 \\
		\hline

	\end{array}
	\]
	
\textbf{Karnaughova mapa}
\begin{center}
	\begin{karnaugh-map}[4][2][1][$BA$][$C$]
		\minterms{0, 1, 2, 4}
		\maxterms{3, 5, 6, 7}
		
		\implicant{0}{4}
		\implicant{0}{1}
		\implicantedge{0}{0}{2}{2}
		
	\end{karnaugh-map}
\end{center}

\textbf{Přepis}\\
	\[x = \neg a \cdot \neg b + \neg a \cdot \neg c + \neg b \cdot \neg c\]
	
\textbf{Schéma}\\

\begin{circuitikz}[scale=0.9, transform shape]
	
	% NOT hradla
	\node[not port] (nota) at (2, 6)  {};
	\node[not port] (notb) at (2, 3)  {};
	\node[not port] (notc) at (2, 0)  {};
	
	% Vstupy
	\draw (nota.in) -- ++(-1.5,0) node[left] {$a$};
	\draw (notb.in) -- ++(-1.5,0) node[left] {$b$};
	\draw (notc.in) -- ++(-1.5,0) node[left] {$c$};
	
	% AND hradla
	\node[and port] (and1) at (6.5, 5)  {};
	\node[and port] (and2) at (6.5, 3)  {};
	\node[and port] (and3) at (6.5, 1)  {};
	
	% Spoje NOT a -> AND1(in1), AND2(in1)
	\draw (nota.out) -- ++(0.4,0) coordinate (Ja);
	\draw (Ja) -- (Ja |- and1.in 1) -- (and1.in 1);
	\draw (Ja) -- (Ja |- and2.in 1) -- (and2.in 1);
	\filldraw (Ja) circle (1.5pt);
	
	% Spoje NOT b -> AND1(in2), AND3(in1)
	\draw (notb.out) -- ++(0.2,0) coordinate (Jb);
	\draw (Jb) -- (Jb |- and1.in 2) -- (and1.in 2);
	\draw (Jb) -- (Jb |- and3.in 1) -- (and3.in 1);
	\filldraw (Jb) circle (1.5pt);
	
	% Spoje NOT c -> AND2(in2), AND3(in2)
	\draw (notc.out) -- ++(0.0,0) coordinate (Jc);
	\draw (Jc) -- (Jc |- and2.in 2) -- (and2.in 2);
	\draw (Jc) -- (Jc |- and3.in 2) -- (and3.in 2);
	\filldraw (Jc) circle (1.5pt);
	
	% OR hradla
	\node[or port] (or1) at (10.5, 4) {};
	\node[or port] (or2) at (14,   2.5) {};
	
	% AND -> OR1
	\draw (and1.out) -- (and1.out |- or1.in 1) -- (or1.in 1);
	\draw (and2.out) -- (and2.out |- or1.in 2) -- (or1.in 2);
	
	% OR1 + AND3 -> OR2
	\draw (or1.out)  -- (or1.out  |- or2.in 1) -- (or2.in 1);
	\draw (and3.out) -- (and3.out |- or2.in 2) -- (or2.in 2);
	
	% Vystup
	\draw (or2.out) -- ++(1,0) node[right] {$x$};
	
	% Popisky
	\node[font=\small, above=3pt] at (and1.north) {$\lnot a \cdot \lnot b$};
	\node[font=\small, above=3pt] at (and2.north) {$\lnot a \cdot \lnot c$};
	\node[font=\small, below=3pt] at (and3.south) {$\lnot b \cdot \lnot c$};
	
\end{circuitikz}
\end{vysvetlenibox}

% ============================================================

% ============================================================
\end{document}
